\documentclass[11pt,a4paper]{article}

% -----------------------------------------------------------------------------
% Corso Linux -- Syllabus
% Documento LaTeX scritto e mantenuto a mano (nessuna generazione con Pandoc).
% -----------------------------------------------------------------------------
\usepackage[utf8]{inputenc}
\usepackage[T1]{fontenc}
\usepackage[italian]{babel}
\usepackage{lmodern}
\usepackage{microtype}
\usepackage[a4paper,margin=2.1cm,headheight=15pt]{geometry}
\usepackage{enumitem}
\usepackage{booktabs}
\usepackage{tabularx}
\usepackage{array}
\usepackage{xcolor}
\usepackage{tikz}
\usetikzlibrary{arrows.meta,calc,positioning,shapes.geometric}
\usepackage[most]{tcolorbox}
\usepackage{titlesec}
\usepackage{fancyhdr}
\usepackage{hyperref}

% --- Identita visiva ----------------------------------------------------------
\definecolor{LinuxBlue}{HTML}{173F5F}
\definecolor{LinuxTeal}{HTML}{168C8C}
\definecolor{LinuxGold}{HTML}{F2B134}
\definecolor{LinuxRed}{HTML}{C94C4C}
\definecolor{LinuxInk}{HTML}{263238}
\definecolor{LinuxPaper}{HTML}{F4F7F8}
\definecolor{LinuxMuted}{HTML}{607D8B}

\hypersetup{
  colorlinks=true,
  linkcolor=LinuxBlue,
  urlcolor=LinuxTeal,
  pdftitle={Corso Linux -- Livello ITS},
  pdfauthor={Giuseppe Capasso},
  pdfsubject={Syllabus del corso Introduzione a Linux}
}

\setlength{\parindent}{0pt}
\setlength{\parskip}{0.55em}
\setlist[itemize]{leftmargin=1.35em,itemsep=0.22em,topsep=0.25em}
\setlist[enumerate]{leftmargin=1.55em,itemsep=0.22em,topsep=0.25em}
\renewcommand{\arraystretch}{1.2}

\titleformat{\section}
  {\Large\bfseries\color{LinuxBlue}}
  {\thesection}{0.65em}{}
\titleformat{\subsection}
  {\large\bfseries\color{LinuxTeal}}
  {\thesubsection}{0.55em}{}
\titlespacing*{\section}{0pt}{1.5em}{0.55em}
\titlespacing*{\subsection}{0pt}{1.1em}{0.35em}

\pagestyle{fancy}
\fancyhf{}
\fancyhead[L]{\small\color{LinuxMuted}Introduzione a Linux}
\fancyhead[R]{\small\color{LinuxMuted}Syllabus}
\fancyfoot[L]{\small\color{LinuxMuted}Giuseppe Capasso}
\fancyfoot[R]{\small\color{LinuxMuted}\thepage}
\renewcommand{\headrulewidth}{0.4pt}
\renewcommand{\footrulewidth}{0pt}

% --- Componenti riutilizzabili -----------------------------------------------
\newtcolorbox{infobox}[1][]{
  enhanced,
  colback=LinuxPaper,
  colframe=LinuxTeal,
  boxrule=0.7pt,
  arc=2mm,
  left=3mm,right=3mm,top=2.2mm,bottom=2.2mm,
  #1
}

\newtcolorbox{lessonbox}[2]{
  enhanced,
  breakable,
  colback=white,
  colframe=LinuxBlue,
  colbacktitle=LinuxBlue,
  coltitle=white,
  fonttitle=\bfseries,
  title={#1\hfill\normalfont\small #2},
  boxrule=0.8pt,
  arc=2mm,
  left=3mm,right=3mm,top=2.5mm,bottom=2.5mm,
  before skip=0.8em,
  after skip=0.8em
}

\newcommand{\term}[1]{\texttt{#1}}
\newcommand{\meta}[2]{\textcolor{LinuxMuted}{\small\textbf{#1}}\\[-0.1em]#2}

\begin{document}

% --- Copertina ----------------------------------------------------------------
\hypersetup{pageanchor=false}
\begin{titlepage}
  \thispagestyle{empty}
  \begin{tikzpicture}[remember picture,overlay]
    \fill[LinuxBlue]
      (current page.north west) rectangle
      ([yshift=-8.2cm]current page.north east);
    \fill[LinuxTeal]
      ([yshift=-8.2cm]current page.north west) rectangle
      ([yshift=-8.55cm]current page.north east);
  \end{tikzpicture}

  \vspace*{1.25cm}
  {\color{white}\large SISTEMI OPERATIVI E AMMINISTRAZIONE}\par
  \vspace{0.55cm}
  {\color{white}\fontsize{34}{39}\selectfont\bfseries
    Introduzione a Linux\par}
  \vspace{0.35cm}
  {\color{white}\LARGE Corso di livello ITS}\par

  \vspace{4.5cm}
  {\Large\bfseries Syllabus del corso}\par
  \vspace{0.45cm}
  {\color{LinuxMuted}\large
    Dalla shell all'amministrazione di un piccolo server Linux}\par

  \vfill
  \begin{infobox}
    \begin{tabularx}{\textwidth}{>{\bfseries\color{LinuxBlue}}l X
                                  >{\bfseries\color{LinuxBlue}}l X}
      Durata & 30 ore & Struttura & 6 lezioni da 5 ore \\
      Destinatari & Studenti ITS, area ICT & Docente & Giuseppe Capasso \\
      Modalita & Teoria essenziale e laboratorio & Livello & Introduttivo \\
    \end{tabularx}
  \end{infobox}

  \vspace{0.8cm}
  {\small\color{LinuxMuted}
    Prerequisiti: conoscenze di base sui sistemi operativi e sui concetti di rete.}
\end{titlepage}

\hypersetup{pageanchor=true}
\pagenumbering{roman}
\tableofcontents
\clearpage
\pagenumbering{arabic}

% --- Presentazione ------------------------------------------------------------
\section{Presentazione del corso}

Il corso introduce Linux attraverso un percorso progressivo e orientato alla
pratica. Ogni lezione collega un nucleo teorico a un laboratorio svolto su
macchine virtuali Ubuntu. Al termine, lo studente sapra utilizzare la shell,
gestire un sistema di base e pubblicare un semplice servizio in rete adottando
misure di sicurezza essenziali.

\subsection{Obiettivi formativi}

Al termine del percorso lo studente sara in grado di:
\begin{itemize}
  \item usare Linux in autonomia e orientarsi nel filesystem;
  \item amministrare utenti, gruppi, file, permessi e processi;
  \item installare, aggiornare e rimuovere software con APT;
  \item configurare storage, mount persistenti e semplici procedure di backup;
  \item diagnosticare la connettivita e configurare accessi remoti;
  \item installare e gestire servizi di sistema;
  \item applicare misure di sicurezza di base e documentare il proprio lavoro.
\end{itemize}

\subsection{Metodo e organizzazione}

La didattica alterna spiegazioni brevi, dimostrazioni del docente, esercitazioni
guidate e problemi da risolvere individualmente o in gruppi di 2--3 persone.
Gli esercizi privilegiano l'interpretazione dell'output e la verifica del
risultato: non e sufficiente copiare un comando, occorre comprenderne effetto e
contesto.

% --- Roadmap TikZ -------------------------------------------------------------
\section{Roadmap del percorso}

La sequenza segue le dipendenze naturali dell'amministrazione Linux: prima si
impara a muoversi nel sistema, poi a gestirlo, collegarlo alla rete e infine a
esporre servizi in modo controllato.

\begin{center}
\begin{tikzpicture}[
  node distance=7mm and 8mm,
  every node/.style={font=\small},
  roadmap/.style={
    rectangle,rounded corners=2mm,minimum width=4.5cm,minimum height=1.55cm,
    text width=3.85cm,align=left,draw=LinuxBlue,line width=0.8pt,fill=white
  },
  num/.style={circle,fill=LinuxGold,draw=none,text=LinuxInk,font=\bfseries,
              minimum size=7mm,inner sep=0pt},
  flow/.style={-{Stealth[length=2.5mm]},line width=1pt,draw=LinuxTeal}
]
  \node[roadmap] (l1) {\textbf{Ambiente e shell}\\VM, distribuzioni, file e directory};
  \node[num,anchor=north east] at ([xshift=-2mm,yshift=-2mm]l1.north east) {1};
  \node[roadmap,right=of l1] (l2) {\textbf{Amministrazione}\\FHS, utenti, permessi e APT};
  \node[num,anchor=north east] at ([xshift=-2mm,yshift=-2mm]l2.north east) {2};
  \node[roadmap,below=of l2] (l3) {\textbf{Storage e backup}\\Partizioni, filesystem, mount, copie};
  \node[num,anchor=north east] at ([xshift=-2mm,yshift=-2mm]l3.north east) {3};
  \node[roadmap,left=of l3] (l4) {\textbf{Networking}\\IP, routing, DNS, SSH e trasferimenti};
  \node[num,anchor=north east] at ([xshift=-2mm,yshift=-2mm]l4.north east) {4};
  \node[roadmap,below=of l4] (l5) {\textbf{Servizi e sicurezza}\\systemd, Apache, firewall};
  \node[num,anchor=north east] at ([xshift=-2mm,yshift=-2mm]l5.north east) {5};
  \node[roadmap,right=of l5] (l6) {\textbf{Consolidamento}\\Progetto di gruppo e prova finale};
  \node[num,anchor=north east] at ([xshift=-2mm,yshift=-2mm]l6.north east) {6};

  \draw[flow] (l1) -- (l2);
  \draw[flow] (l2) -- (l3);
  \draw[flow] (l3) -- (l4);
  \draw[flow] (l4) -- (l5);
  \draw[flow] (l5) -- (l6);
\end{tikzpicture}
\end{center}

\begin{infobox}
  \textbf{Distribuzione indicativa delle 30 ore:}
  circa 12 ore di spiegazione e dimostrazione, 15 ore di laboratorio e 3 ore
  dedicate a verifica, restituzione e consolidamento. Il docente puo adattare
  il ritmo in base all'esperienza della classe.
\end{infobox}

% --- Programma ---------------------------------------------------------------
\section{Programma delle lezioni}

\begin{lessonbox}{Lezione 1 --- Installazione e ambiente di lavoro}{5 ore}
  \textbf{Risultato atteso.} Lo studente prepara una VM Ubuntu, riconosce gli
  elementi di una distribuzione e svolge operazioni essenziali nella shell.

  \begin{tabularx}{\textwidth}{>{\bfseries\color{LinuxTeal}}p{3cm} X}
    Concetti & Linux e GNU/Linux; distribuzioni; ambiti di utilizzo;
               virtualizzazione; terminale e shell. \\
    Comandi & \term{ls}, \term{pwd}, \term{cd}, \term{mkdir}, \term{touch},
              \term{cp}, \term{mv}, \term{rm}, \term{ls -l}, \term{chmod}. \\
    Laboratorio & Installazione o verifica della VM; navigazione del filesystem;
                  creazione di una struttura personale di file e directory;
                  prime modifiche ai permessi. \\
  \end{tabularx}
\end{lessonbox}

\begin{lessonbox}{Lezione 2 --- Amministrazione di base e APT}{5 ore}
  \textbf{Risultato atteso.} Lo studente gestisce account locali, applica
  permessi coerenti e amministra software dai repository.

  \begin{tabularx}{\textwidth}{>{\bfseries\color{LinuxTeal}}p{3cm} X}
    Concetti & Struttura FHS; directory principali; tipi di file; utenti,
               gruppi e privilegi; package manager e repository. \\
    Comandi & \term{adduser}, \term{deluser}, \term{groups}, \term{sudo},
              \term{chmod}, \term{chown}, \term{chgrp}; \term{apt update},
              \term{apt upgrade}, \term{apt install}, \term{apt remove},
              \term{apt search}. \\
    Laboratorio & Creazione di utenti e gruppi; verifica dei permessi;
                  ricerca, installazione e rimozione controllata di pacchetti. \\
  \end{tabularx}
\end{lessonbox}

\begin{lessonbox}{Lezione 3 --- Storage e backup}{5 ore}
  \textbf{Risultato atteso.} Lo studente prepara un disco aggiuntivo, lo monta
  in modo persistente ed esegue backup e ripristino di una directory.

  \begin{tabularx}{\textwidth}{>{\bfseries\color{LinuxTeal}}p{3cm} X}
    Concetti & HDD e SSD; device in \term{/dev}; MBR e GPT; filesystem ext4;
               mount temporaneo e persistente; backup completo e incrementale. \\
    Strumenti & \term{lsblk}, \term{fdisk}/\term{cfdisk}, \term{mkfs},
                \term{/etc/fstab}, \term{rsync}, \term{tar}. \\
    Laboratorio & Aggiunta di un disco virtuale; partizionamento e formattazione;
                  mount permanente; backup e ripristino di file selezionati. \\
  \end{tabularx}

\end{lessonbox}

\begin{lessonbox}{Lezione 4 --- Networking Linux}{5 ore}
  \textbf{Risultato atteso.} Lo studente interpreta la configurazione di rete,
  individua problemi di connettivita e trasferisce file tramite SSH.

  \begin{tabularx}{\textwidth}{>{\bfseries\color{LinuxTeal}}p{3cm} X}
    Concetti & DHCP e configurazione statica; indirizzo, rete e gateway; routing;
               risoluzione dei nomi; accesso remoto sicuro. \\
    Strumenti & \term{ip a}, \term{ip r}, \term{ping}, \term{traceroute},
                \term{ss};\newline \term{/etc/hosts}, \term{resolv.conf},
                \term{nslookup}, \term{ssh}, \term{scp}. \\
    Laboratorio & Collegamento tra VM; diagnosi con ping e route; associazione di
                  nomi locali; accesso remoto e copia di un file. \\
  \end{tabularx}
\end{lessonbox}

\begin{lessonbox}{Lezione 5 --- Servizi e sicurezza di base}{5 ore}
  \textbf{Risultato atteso.} Lo studente pubblica una pagina web, gestisce il
  relativo servizio e limita l'accesso alle sole porte necessarie.

  \begin{tabularx}{\textwidth}{>{\bfseries\color{LinuxTeal}}p{3cm} X}
    Concetti & Processi in background e demoni; avvio dei servizi; web server;
               porte, superficie esposta e regole firewall. \\
    Strumenti & \term{systemctl start|stop|enable|status}; Apache;
                \term{/var/www/html}; \term{ufw allow|deny|status}. \\
    Laboratorio & Installazione di Apache; modifica di \term{index.html};
                  apertura controllata della porta 80; test da una seconda VM. \\
  \end{tabularx}
\end{lessonbox}

\begin{lessonbox}{Lezione 6 --- Consolidamento e lavoro di gruppo}{5 ore}
  \textbf{Risultato atteso.} Il gruppo integra le competenze del corso in un
  piccolo sistema funzionante e ne dimostra configurazione e sicurezza.

  \begin{tabularx}{\textwidth}{>{\bfseries\color{LinuxTeal}}p{3cm} X}
    Organizzazione & Gruppi di 2--3 studenti, ruoli espliciti, checklist tecnica
                    e breve restituzione finale. \\
    Progetto & Preparare una VM; configurare rete, utenti e permessi; installare
               software; aggiungere storage; pubblicare un sito; consentire solo
               SSH e HTTP; verificare servizio e firewall. \\
    Consegna & Evidenze dei controlli eseguiti, elenco dei comandi significativi
               e spiegazione delle decisioni adottate. \\
  \end{tabularx}
\end{lessonbox}

% --- Competenze ---------------------------------------------------------------
\section{Riepilogo delle competenze}

Al termine delle sei lezioni lo studente dispone delle seguenti competenze
operative di base.

\begin{tabularx}{\textwidth}{>{\bfseries\color{LinuxBlue}}p{3.5cm} X}
  Shell e filesystem & Navigare tra directory; creare, copiare, spostare e
                       rimuovere file; leggere percorsi e permessi. \\
  Utenti e permessi & Gestire account e gruppi; assegnare proprietari e permessi;
                      usare \term{sudo} in modo consapevole. \\
  Gestione software & Cercare, installare, aggiornare e rimuovere pacchetti dai
                      repository con APT. \\
  Storage e backup & Riconoscere dischi e partizioni; creare un filesystem;
                     configurare un mount persistente; eseguire backup e restore. \\
  Rete & Leggere indirizzi e route; verificare connettivita e DNS; usare SSH e
         trasferire file tra due sistemi. \\
  Servizi e sicurezza & Gestire un servizio con \term{systemctl}; pubblicare una
                        pagina con Apache; configurare regole essenziali con UFW. \\
\end{tabularx}

\begin{infobox}[title={Profilo in uscita}]
  Le competenze acquisite costituiscono una base per attivita junior di
  \textbf{help desk}, \textbf{supporto IT} e \textbf{amministrazione di sistema}.
\end{infobox}

\end{document}
